\documentclass{article}
\usepackage[UTF8]{ctex}
\usepackage{amsmath,amssymb,amsfonts}
\usepackage{enumitem}
\usepackage{geometry}
\geometry{a4paper, margin=1in}
% === 错题本排版优化（不参与题目标识） ===
\setlength{\parindent}{0pt}
\setlength{\parskip}{0.35em plus 0.10em minus 0.08em}
\setlength{\emergencystretch}{3em}
\setlist{itemsep=0.25em, topsep=0.40em, parsep=0pt}
\allowdisplaybreaks[2]
% === 错题本排版优化结束 ===
\begin{document}
\section*{安徽大学2025年高等代数真题}
\section*{104．安徽大学2025年高等代数试卷}
一．填空题．每题5分，共30分．\\

\subsection{题目1}
1．多项式 \(x^{6}+1\) 在 \(\mathbb{R}\) 上的不可约分解为 \(\_\_\_\_\)．

\subsection{题目2}
2．已知 5 阶方阵满足 \(A^{2}-3 A+2 I=O\) ，其中 \(I\) 为单位阵，若迹 \(\operatorname{tr}(A-I)=2\) ，则 \(A\) 的特征值1的几何重数为 \(\_\_\_\_\)。

\subsection{题目3}
3．若线性方程组 \(\left\{\begin{array}{l}\lambda x_{1}+x_{2}+x_{3}=1 ; \\ x_{1}+\lambda x_{2}+x_{3}=\lambda ; \\ x_{1}+x_{2}+\lambda x_{3}=\lambda^{2} .\end{array}\right.\) 无解，则 \(\lambda\) 的值为 \(\_\_\_\_\)．\\
\(\_\_\_\_\)

\subsection{题目4}
4．设 3 阶方阵 \(A\) 恰好有一个 Jordan 块，则 \(A\) 有 \(\_\_\_\_\)个不同的不变因子．

\subsection{题目5}
5．向量组（1，1，－1，－1），（1，1，0，0），（0，0，－1，－1），（2，2，－1，－1），（1，0，－1，－1）在四维行向量空间（取\(\mathbb{R}^{4}\) 标准内积）中生成的子空间的正交补空间的维数为 \(\_\_\_\_\)．

\subsection{题目6}
6．设 \(n\) 阶矩阵 \(A\) 满足秩 \((A+I)+\) 秩 \((A-I)=n\) ，且 \(A \neq I\) ，其中 \(I\) 为单位矩阵，则 \(A\) 必有特征值 \(\_\_\_\_\)．

二．辨析题．每题 5 分，回答正确或错误 2 分，简单说明原因 3 分．共 20 分．\\

\subsection{题目7}
7．设 \(U, V\) 为数域 \(P\) 上的线性空间，\(\varphi: U \rightarrow V\) 为线性映射，则商空间 \(V / \operatorname{Im} \varphi\) 与 \(\operatorname{Ker} \varphi\) 同构。

\subsection{题目8}
8．设 \(V_{1}, V_{2}\) 分别是数域 \(P\) 上齐次线性方程组 \(x_{1}=x_{2}=\cdots=x_{n}\) 与 \(x_{1}+x_{2}+\cdots+x_{n}=0\) 的解空间，则 \(P^{n}=V_{1} \oplus V_{2}\) 。

\subsection{题目9}
9．设 \(A\) 为 \(m \times n\) 矩阵，\(A^{\mathrm{T}}\) 为 \(A\) 的转置，则秩 \(\left(A^{\mathrm{T}} A\right)=\) 秩 \((A)\) 。

\subsection{题目10}
10．设 \(A\) 为正交阵，且行列式 \(|A|=-1\) ，则－1必为 \(A\) 的特征值。
\end{document}
